% !TEX root = ../main.tex
\chapter{Regularized Pricing and Risk Models}\label{ch:regularized}
\begin{paracol}{2}

\begin{sources}
\begin{tabularx}{\linewidth}{@{}l l L@{}}
\toprule
\textbf{Lecture} & \textbf{Speaker} & \textbf{Content}\\
\midrule
2013 L10 (1 h 30) & I.~Masyukov (Morgan Stanley) & Guest lecture. Bonds, yield, duration, convexity; interest-rate swaps and the swap rate as an average of forward rates; building a yield curve with splines; bond spread over the swap curve; the cost of hedging a swap book; PCA hedging and its instabilities; a hedging model regularised by smoothness of the forward curve; the pricing/calibration engine and the HJM model; calibration of a two-dimensional volatility surface as an ill-posed inverse problem; SVD analysis, Tikhonov regularisation, truncated SVD; questions on splines, relative-value trading and bond curves.\\
\bottomrule
\end{tabularx}

\smallskip
\textbf{Files.} Slides \file{MIT18_S096F13_lecnote10.pdf} (51 pages, no incremental builds).
\textbf{Problem sets.} None. The exercises (each after the section it practises) are all ``(not from the course)''.
\textbf{Not published.} The talk has no 2024 counterpart. The slides give the numerical tables and figures
of the swap-book example, but no data or code; the toy computations of this chapter (a swap-curve
Jacobian, a one-dimensional calibration) are my own and are marked as such.
\end{sources}

\section*{Motivation}
Most quantities a bank computes every day are the output of an \emph{inverse problem}. A model has
parameters (a yield curve, a volatility surface); a handful of liquid instruments have observed prices;
the parameters must be chosen so that the model reproduces the prices, and only then can it be used to value and
hedge a book of thousands of trades that are \emph{not} among the calibration instruments. The forward map
``parameters $\mapsto$ prices'' is smooth and easy; its inverse is not. Small errors in the prices, or small
differences between instruments, are amplified into wild oscillations of the parameters, and the model that
``fits perfectly'' prices the rest of the book badly. The cure, the subject of this guest lecture, is
\emph{regularisation}: add to the fit a penalty that says what a sensible solution looks like, and give up an exact
fit in exchange for stability.

The chapter follows the lecture. \Cref{ch20:sec-bonds} recalls bonds, swaps and the yield curve. \Cref{ch20:sec-curve}
shows on a toy how calibrating a curve exactly already amplifies a 1~bp change into a ripple of several
basis points. \Cref{ch20:sec-hedge} treats the risk side: hedging with a few liquid swaps by PCA and by a smoothness-regularised
alternative. \Cref{ch20:sec-inverse} develops the mathematics behind both, singular values, Tikhonov regularisation and
truncated SVD, which is the natural continuation of the shrinkage methods of \cref{ch05:sec-shrink}; \cref{ch20:sec-lambda}
discusses how to choose the regularisation parameter (not covered in class). \Cref{ch20:sec-vol} applies everything
to the calibration of a volatility surface for the Heath--Jarrow--Morton model (\cref{ch:hjm}). The linear algebra needed is
the SVD of \cref{sec:svd}.

% ==================================================================
\section{Bonds, swaps and the yield curve}\label{ch20:sec-bonds}

\subsection{Bond price, yield, duration and convexity}\label{ch20:sec-yield}
A fixed-rate bond pays cash flows $C_1,\dots,C_N$ (coupons, plus the face value in $C_N$) at times $t_1<\dots<t_N$
\ts{13}{10}{2:11}. Cash paid later is worth less, so each cash flow is multiplied by a \emph{discount factor}
$\Lambda_i\in(0,1]$ that decreases with $t_i$ \ts{13}{10}{5:16}, and the fair value is
$P=\sum_{i=1}^N\Lambda_iC_i$ \ts{13}{10}{6:20}. For a bond the price $P$ is not something to be computed: it is the
result of a liquid market and is known \ts{13}{10}{7:20}. What is needed is a \emph{model for discounting}. The simplest one uses a single
parameter, the \emph{yield to maturity} $y$ (continuous compounding, one rate for all dates)
\ts{13}{10}{8:27}:
\begin{equation}\label{ch20:eq-yield}
  \Lambda_i=e^{-yt_i},\qquad P(y)=\sum_{i=1}^Ne^{-yt_i}C_i .
\end{equation}
With $C_i>0$ the map $y\mapsto P(y)$ is strictly decreasing, so price and yield determine each other one-to-one
\ts{13}{10}{9:31}. The yield is therefore not a traded quantity but a \emph{convention} to summarise price and
cash flows \ts{13}{10}{10:34}. Differentiating,
\[
  \frac{\dd P}{\dd y}=-\sum_it_ie^{-yt_i}C_i,\qquad
  D=-\frac1P\frac{\dd P}{\dd y}=\sum_{i=1}^Nt_i\,\pi_i,\quad \pi_i=\frac{e^{-yt_i}C_i}{P},\ \ \sum_i\pi_i=1,
\]
\[
  \mathcal C=\frac1P\frac{\dd^2P}{\dd y^2}=\sum_it_i^2\pi_i>0 .
\]
The \emph{duration} $D$ is the present-value-weighted average time of the cash flows (``weighted time'')
\ts{13}{10}{11:40}\ts{13}{10}{12:41}; normalising by $P$ makes it independent of how many bonds are held. A zero-coupon
bond has one weight, so $D=T$; a coupon bond has $D<t_N$, because every $t_i\le t_N$ and at least one is strictly smaller
\ts{13}{10}{13:47}. Convexity is a weighted mean of $t_i^2$, hence positive. These are the definitions of
\cref{sec:duration} in the continuous-compounding convention (\cref{prop:duration}), and \cref{eq:durconv} is the Taylor expansion
$\Delta P/P\approx-D\Delta y+\frac12\mathcal C(\Delta y)^2$.

\begin{coursenote}[Sign and letters in the slides]
Slide~8 defines ``duration'' as $d=\frac1P\frac{\partial P}{\partial y}$, which is \emph{negative} (the speaker says so:
the duration is always negative); here, as in \cref{sec:duration}, $D=-d>0$. Slide~9 names the second derivative $c=\partial^2P/\partial y^2$
and then writes its explicit form as ``$d=\sum t_i^2e^{-yt_i}C_i$'': the letter should be $c$, and this is the
\emph{unnormalised} second derivative, $\mathcal C=c/P$ in the notation used here.
\end{coursenote}

\begin{example}[Duration explains most of the P\&L, convexity the rest (my numbers)]\label{ch20:ex-bond}
Take $N=5$ annual payments $C=(4,4,4,4,104)$ and $y=4\%$. Then $P=99.640$, $D=4.629$ years, $\mathcal C=22.42$. If $y$ rises by $50$~bp the exact
price change is $-2.2868\%$; the duration term alone gives $-2.3146\%$ and adding convexity $-2.2865\%$. For $+200$~bp the figures are
$-8.824\%$ exactly, $-9.258\%$ with duration only and $-8.810\%$ with convexity. The trader who
explains the day's loss with first derivatives only \ts{13}{10}{15:56} is left with an ``unexplained'' term
of order $\frac12\mathcal C(\Delta y)^2$.
\end{example}

All of this assumes one $y$ for all dates, hence that the curve moves in \emph{parallel}. That was an acceptable
approximation before the crisis, but with a steep curve, where short rates sit near zero and long rates are expected to
rise, it is not \ts{13}{10}{14:55}. \Cref{ch20:sec-bondspread} replaces $y$ by a full curve.

\subsection{Interest-rate swaps}\label{ch20:sec-swaps}
In a swap one party pays a fixed rate $C$ and receives a floating rate (three-month LIBOR in the lecture, fixed on the
day and known in advance for the accrual period) \ts{13}{10}{17:02}. With payment dates $T_1<\dots<T_n$, accrual
fractions $\delta_i$ (day-count), discount factors $\Lambda_i$ and forward rates $r_i$ for the period ending at $T_i$,
\[
  PV_{\rm fixed}=\sum_iC\delta_i\Lambda_i=C\sum_iw_i,\qquad PV_{\rm float}=\sum_ir_i\delta_i\Lambda_i=\sum_ir_iw_i,\qquad w_i=\delta_i\Lambda_i,
\]
\ts{13}{10}{18:14}. A swap costs nothing to enter, so the fixed rate is set so that the two legs have the same value
\ts{13}{10}{19:21}:
\begin{equation}\label{ch20:eq-swaprate}
  C=\frac{\sum_ir_iw_i}{\sum_iw_i}\qquad(\text{same payment frequency on both legs}),
\end{equation}
the \emph{swap rate}, a weighted average of forward rates, which is what traders quote \ts{13}{10}{20:27}. If the forward
rates are defined by $\Lambda_{i-1}=(1+r_i\delta_i)\Lambda_i$ (the single-curve convention of \cref{sec:forward}, $\Lambda_0=1$), the floating leg telescopes,
$PV_{\rm float}=\sum_i(\Lambda_{i-1}-\Lambda_i)=1-\Lambda_n$, and $C=(1-\Lambda_n)/\sum_i\delta_i\Lambda_i$ (my remark, not on the slides). The fixed leg is a bond coupon stream, which is why a swap can be
hedged with a bond \ts{13}{10}{20:27}, and why swap risk is measured in \emph{buckets} of par rates: the trader wants to know
the P\&L if the 5-year swap rate moves by 1~bp.

\subsection{Building a curve}\label{ch20:sec-build}
A liquid market gives swap rates for a few maturities, whereas a book of hundreds of thousands of swaps needs a discount factor for every day up to 30--40 years
\ts{13}{10}{22:35}. Constructing the curve means \ts{13}{10}{23:39}:
\begin{enumerate}[1.]
  \item choose the input instruments;
  \item choose
      which quantity is interpolated (daily forward rates, zero rates, discount factors) and the interpolation family
      (piecewise-constant, linear, tension spline, cubic spline) with its \emph{knots};
  \item \emph{calibrate}: solve for the
      spline parameters so that every input instrument is repriced at par (PV of the swap equal to zero)
      \ts{13}{10}{24:42}\ts{13}{10}{28:58}. The lecture's example uses 15 swaps from 1 to 30 years with quotes from $0.33\%$ to $2.67\%$
      (the market of a year earlier; \cref{ch20:tab-risk}); the resulting curve of three-month forwards is far from flat, steep in the first five years
      and with a hump in the 20-year region \ts{13}{10}{30:01}.
\end{enumerate}

\paragraph{Splines and B-splines.}
A \emph{cubic spline} is a function that is a cubic polynomial on each interval between knots and has two continuous derivatives at every knot \ts{13}{10}{25:48}.
The knots in the lecture's picture are $1,10,20,40,80,160,240$ (quarters), denser where instruments are dense
\ts{13}{10}{1:22:09}. All splines with given knots form a vector space, which has a convenient basis of bell-shaped,
locally supported functions \ts{13}{10}{26:49}\ts{13}{10}{27:51}.

\begin{definition}[B-splines, Cox--de Boor recursion]\label{ch20:def-bspline}
For a non-decreasing knot sequence $\tau_0\le\tau_1\le\cdots$ let $B_{i,0}(t)=1$ if $\tau_i\le t<\tau_{i+1}$ and $0$ otherwise, and
\[
  B_{i,k}(t)=\frac{t-\tau_i}{\tau_{i+k}-\tau_i}B_{i,k-1}(t)+\frac{\tau_{i+k+1}-t}{\tau_{i+k+1}-\tau_{i+1}}B_{i+1,k-1}(t)
\]
(a term with zero denominator is dropped) \ts{13}{10}{1:08:13}. Degree $k=3$ gives cubic B-splines.
\end{definition}

\begin{proposition}[Properties of B-splines]\label{ch20:prop-bspline}
$B_{i,k}\ge0$; it vanishes outside $[\tau_i,\tau_{i+k+1}]$; $\sum_iB_{i,k}(t)=1$ on the interior of the knot range; for simple knots
$B_{i,k}\in C^{k-1}$. Hence every $f=\sum_ic_iB_{i,k}$ is a spline of degree $k$ and its values lie between the smallest and the largest $c_i$.
\end{proposition}
\begin{proof}
Non-negativity and support follow by induction from the recursion, whose coefficients are non-negative on the supports involved. For the partition of unity
shift the index in the second term of $\sum_iB_{i,k}$:
\[
  \sum_iB_{i,k}=\sum_iB_{i,k-1}\Bigl[\frac{t-\tau_i}{\tau_{i+k}-\tau_i}+\frac{\tau_{i+k}-t}{\tau_{i+k}-\tau_i}\Bigr]=\sum_iB_{i,k-1},
\]
and the case $k=0$ is $1$. Smoothness: at a simple knot each recursion step raises the differentiability by one, starting from a jump. (Checked numerically: the maximal deviation from $1$ is $4\times10^{-16}$.)
\end{proof}
Local support makes the calibration matrix banded and the response to a change of one parameter local; the convex-combination
property makes the parameters interpretable as ``values of the curve''. This is why the same basis is used for the surface, in two dimensions,
in \cref{ch20:sec-vol}.

\subsection{Bond spread over the curve}\label{ch20:sec-bondspread}
Once a swap curve $\Lambda_i$ exists, a bond can be priced off it with one extra parameter, the \emph{bond spread} $s$,
determined so that the model reproduces the (liquid, known) bond price \ts{13}{10}{31:01}:
\begin{equation}\label{ch20:eq-spread}
  P=\sum_{i=1}^Ne^{-st_i}\Lambda_iC_i .
\end{equation}
Compared with $y$ the correction is small: a curve level of $3\%$ and a spread of a few basis points, about a hundred times smaller
\ts{13}{10}{32:10}. The gain is consistency: sensitivities of the swap book to the curve inputs are inherited by the bond, and one is no longer bound to
the single ``parallel'' mode of duration. The spread contains credit or liquidity components of the bond, and even a
swap-versus-Treasury spread is traded; for a 10-year maturity the most liquid instrument is the bond, then, surprisingly, the spread, and only then the swap
itself, so the curve is built from bond yield plus spread rather than from the raw 10-year swap quote
\ts{13}{10}{33:10}\ts{13}{10}{34:11}\ts{13}{10}{35:12}.

\begin{finance}[Why the curve is built on swaps and not on bonds]
Asked why one does not invert the bond curve instead, the speaker gave two reasons \ts{13}{10}{1:27:24}\ts{13}{10}{1:28:29}. A swap
starting today always exists, and tomorrow's swap starts tomorrow, whereas bonds have a fixed maturity: the newest issue (\emph{on-the-run}) is by far the most liquid
and older ones (\emph{off-the-run}) trade at different levels, so there is no continuous family of maturities and a bond curve can only be a least-squares fit, not an exact one.
The switch of liquidity from one issue to the next creates ``roll effects'' that make the curve unstable. So the swap desk builds the swap curve and \emph{projects} bonds onto it, by the spread
\eqref{ch20:eq-spread}, since bonds are its preferred hedging instruments.
\end{finance}

% ==================================================================
\section{Calibrating a curve exactly is already ill-conditioned}\label{ch20:sec-curve}

The lecture shifts one input of the calibrated curve, the 9-year swap rate, by $1$~bp, keeping all the other 14 quotes fixed, and plots the response of the
three-month forward rates: it is a ripple, zero before about 5 years and after about 20 years, with a swing of about $\pm14$~bp around the 9-year point
\ts{13}{10}{36:16}\ts{13}{10}{37:22}. The curve is exact at every input, so it \emph{must} move in the neighbourhood of the shifted input (the 8- and 10-year par
rates are pinned), and it does so by an oscillation much larger than the shift. The speaker's comment: this is an ill-posed problem, where a
small change of the input causes a large change of the output.

\paragraph{A toy where everything is explicit (not in class).}\label{ch20:par-toy}
Let the instantaneous forward rate be constant on each interval $(T_{i-1},T_i]$, $T_0=0$, between consecutive input maturities, with value $f_i$, and let the par rate of the swap
maturing at $T_i$ be (ignoring discounting) the average forward rate, $S_i=\frac1{T_i}\sum_{k\le i}(T_k-T_{k-1})f_k$, i.e. $S=Mf$ with $M$ lower triangular
and invertible. Inverting,
\[
  f_i=\frac{T_iS_i-T_{i-1}S_{i-1}}{T_i-T_{i-1}},\qquad\text{so}\qquad J=M^{-1}:\quad \Delta f_i=\frac{T_i\,\Delta S_i-T_{i-1}\,\Delta S_{i-1}}{T_i-T_{i-1}} .
\]
A shift $\Delta S_j=1$~bp moves $f_j$ by $+T_j/(T_j-T_{j-1})$ and $f_{j+1}$ by $-T_j/(T_{j+1}-T_j)$: for the 9-year input, $+9$~bp on $(8,9]$ and $-9$~bp on
$(9,10]$; for 10 years $+10$ and $-5$; for 25 years $+5$ and $-5$. Each maturity acts as a lever with arm $T_j$ divided by the spacing
of the instruments, and no amount of care in the method removes it, because the data (an average up to $T_j$) constrain only a
difference of averages. For the 15 maturities of the lecture the condition number of $M$ is $23.4$ (my computation). \Cref{ch20:fig-ripple}
draws the two responses.

\begin{figure}[ht]
\centering
\begin{tikzpicture}
\begin{axis}[width=12cm,height=4.6cm,xlabel={forward time $t$ (years)},ylabel={$\Delta f(t)$ (bp)},xmin=0,xmax=30,ymin=-10.5,ymax=10.5,
  axis lines=left,grid=major,grid style={black!12},legend style={at={(0.5,-0.42)},anchor=north,legend columns=2,draw=none},tick label style={font=\small}]
\addplot[thick,thmcol] coordinates {(0,0)(8,0)(8,9)(9,9)(9,-9)(10,-9)(10,0)(30,0)};
\addlegendentry{$+1$~bp on the 9Y par rate}
\addplot[thick,defcol,dashed] coordinates {(0,0)(20,0)(20,5)(25,5)(25,-5)(30,-5)};
\addlegendentry{$+1$~bp on the 25Y par rate}
\end{axis}
\end{tikzpicture}
\caption{Forward-rate response to a $1$~bp shift of one par rate in the piecewise-constant toy (my computation): the curve calibrated exactly answers with a
ripple whose amplitude is the maturity divided by the spacing of instruments. The lecture's spline curve gives a similar but smoother ripple (\ts{13}{10}{36:16}).}\label{ch20:fig-ripple}
\end{figure}

\begin{remark}[Smoother is not automatically better (my computation)]
Replacing the piecewise-constant forwards by a cubic B-spline with 15 coefficients (interior knots by the averaging rule of de Boor) and again calibrating exactly
gives a condition number of about $6\times10^3$ for the calibration matrix, and the response to the same $1$~bp shift of the 9-year rate has amplitude of
between $10$ and $30$~bp at $9$--$10$~years, about $70$~bp at $12$, $150$~bp at $16$ and $200$~bp at $20$ years. The exact calibration of a few averaged quantities to a flexible smooth
curve is genuinely unstable in the sparse long end; the lecture's curve (ripple confined to $\pm14$~bp) shows that the design of the interpolation matters.
\end{remark}

\begin{exercise}[(not from the course)]
\leavevmode
\begin{enumerate}[(a)]
  \item For a swap with equally spaced payment dates show $C=(1-\Lambda_n)/\sum_i\delta_i\Lambda_i$ and that $C$ lies between the smallest and the largest forward rate.
  \item In the piecewise-constant toy of \cref{ch20:sec-curve} derive $\Delta f_i$ for a shift of the $j$-th par rate and show that its entries are $T_j/(T_j-T_{j-1})$ and $-T_j/(T_{j+1}-T_j)$; what is the largest absolute value over all $j$ for the 15 maturities of the lecture?
  \item The matrix $M$ is lower triangular: compute its determinant.
\end{enumerate}
\end{exercise}

% ==================================================================
\section{Hedging a swap book: PCA and the smoothness-regularised model}\label{ch20:sec-hedge}

\subsection{The cost of hedging bucket by bucket}\label{ch20:sec-cost}
The \emph{risk} of a book is the vector $\bm r\in\R^n$ of P\&L per basis point for a shift of each of the $n$ input rates (here $n=15$ swap maturities);
to first order the P\&L for a shift $\bm u$ of the inputs is $\bm r\T\bm u$. \Cref{ch20:tab-risk} shows the lecture's example \ts{13}{10}{37:22}. The
net risk is tiny ($2{,}266$ in total), but the risk in each bucket is large and of both signs. To be insensitive to \emph{every} bucket the trader must trade
$\pm r_i$ swaps in each of them, paying half the bid--offer of each instrument; the liquid maturities (yellow in the slide: 1, 2, 5, 10, 30 years) cost $0.10$~bp, the
others $0.25$~bp. The charge of a bucket is $|r_i|\times$ bid--offer, and the total is $\sum_i|r_i|b_i=3{,}617{,}629$
\ts{13}{10}{38:25}\ts{13}{10}{39:27}. The gross exposure, not the net one, is what costs money, so traders never hedge every bucket: they hedge a seven-year risk
with the cheaper liquid points around it \ts{13}{10}{40:30}. The question is how to do this systematically.

\begin{table}[ht]
\centering\small
\begin{tabular}{lrrrr}
\toprule
Swap & quote (\%) & raw risk & bid--offer (bp) & charge\\
\midrule
1Y & 0.33 & $-200{,}000$ & 0.10 & 20,000\\
2Y & 0.39 & 1,330,000 & 0.10 & 133,000\\
3Y & 0.49 & $-200{,}000$ & 0.25 & 50,000\\
4Y & 0.64 & 1,200,000 & 0.25 & 300,000\\
5Y & 0.86 & $-722{,}450$ & 0.10 & 72,245\\
6Y & 1.09 & $-35{,}255$ & 0.25 & 8,814\\
7Y & 1.29 & $-537{,}430$ & 0.25 & 134,358\\
8Y & 1.48 & $-3{,}850{,}000$ & 0.25 & 962,500\\
9Y & 1.64 & 1,580,000 & 0.25 & 395,000\\
10Y & 1.79 & 288,751 & 0.10 & 28,875\\
12Y & 2.04 & $-401{,}350$ & 0.25 & 100,338\\
15Y & 2.29 & 50,000 & 0.25 & 12,500\\
20Y & 2.50 & 4,000,000 & 0.25 & 1,000,000\\
25Y & 2.60 & $-1{,}000{,}000$ & 0.25 & 250,000\\
30Y & 2.67 & $-1{,}500{,}000$ & 0.10 & 150,000\\
\midrule
total & & 2,266 & & 3,617,629\\
\bottomrule
\end{tabular}
\caption{The swap book of the lecture (slide~15; the unit of the raw risk is dollars per basis point, presumably). The charge is $|\text{risk}|\times\text{bid--offer}$;
I checked that the column sums to $3{,}617{,}629$.}\label{ch20:tab-risk}
\end{table}

\subsection{Hedging with a few liquid swaps: the general formulation}\label{ch20:sec-hedgeform}
Let $H\in\R^{n\times k}$ have as columns the risk vectors (in the $n$ buckets) of $k$ liquid hedging instruments; if the same instruments calibrate the curve, a hedge instrument
is sensitive only to its own input, and $H$ is made of columns $\bm e_{i(j)}$ of the identity (slide~19). Take $k=5$: the 1, 2, 5, 10 and 30 year swaps. Let $F\in\R^{n\times m}$ collect market
\emph{scenarios} (moves of the inputs) \ts{13}{10}{41:37}. Holding a hedge position that transfers risk $H\bm x$, the residual risk is $\bm r-H\bm x$ and its P\&L in scenario $\bm f$ is $\bm f\T(\bm r-H\bm x)$, so one
chooses
\begin{equation}\label{ch20:eq-hedgeopt}
  \bm x=\arg\min_{\bm x}\ \bigl\|F\T(\bm r-H\bm x)\bigr\|^2 .
\end{equation}
The vector $\bm x$ (the ``translated risk''; the hedge trades are $-\bm x$) is the equivalent risk of the book in the $k$ liquid swaps.

\begin{coursenote}[Sign convention on slides 16 and 19]
Slide~16 writes the objective as $\|F\T(\bm r+H\bm x)\|^2$ and slides 19--23 use $\bm r-H\bm x$ (and $\bm x=R\T\bm r$). I use the latter throughout; the two differ by $\bm x\to-\bm x$.
\end{coursenote}

\subsection{PCA hedging}\label{ch20:sec-pcahedge}
Take for $F$ the first principal components of historical daily changes of the inputs. If $D=USP\T$ is the SVD of the data matrix of changes (rows are days),
$P$ has as columns the principal components, which are eigenvectors of $D\T D$ (\cref{ch:pca}, \cref{thm:svd}) \ts{13}{10}{42:50}; keeping $k$ of them, one for each hedging instrument, is a low-rank
approximation of the market. In the lecture's data the first component is the ``rates do not move now but move in the future'' mode (steep up to 10 years,
flat afterwards, because the Fed is holding), the second is a tilt, the third more complex \ts{13}{10}{43:57}\ts{13}{10}{45:02}.

With $F=P\in\R^{n\times k}$ the objective \eqref{ch20:eq-hedgeopt} can be made \emph{exactly} zero, since $P\T H\in\R^{k\times k}$ is (generically) invertible \ts{13}{10}{45:02}:
\begin{equation}\label{ch20:eq-pcarisk}
  P\T(\bm r-H\bm x)=\bm 0\ \Longrightarrow\ \bm x=(P\T H)^{-1}P\T\bm r=R\T\bm r,\qquad R=P(H\T P)^{-1}\in\R^{n\times k}.
\end{equation}
The matrix $R$ deserves an interpretation (slide~22, \ts{13}{10}{47:12}): it is the unique combination $R=PY$ of the components with $H\T R=I_k$. Its $j$-th column is the
\emph{scenario} (a shift of all 15 inputs, as a combination of principal components) in which the $j$-th hedging swap moves by $1$~bp and the other four do not move at all; and $x_j=\bm r\T R_{\cdot j}$ is the P\&L of the book in that scenario.
Rows of $R$ are the weights with which the risk of each bucket is assigned to the liquid swaps. This is the matrix of slide~20, whose rows have the unit vectors at the hedging tenors (a swap hedges itself)
and, for example, assigns the 4-year risk mainly to the 2-year and 5-year swaps, with weights $0.60$ and $0.70$. The translated risks and their charges are in \cref{ch20:tab-hedge}
(first block): total charge $442{,}036$ against $3{,}617{,}629$, a reduction of a factor $8.2$.

\begin{coursenote}[``Orders of magnitude'']
The speaker says that the charge becomes ``orders of magnitude smaller'', ``something like $400$'' instead of $3.6$ million \ts{13}{10}{47:12}. The slide gives $442{,}036$ for the PCA model
(and $446{,}717$ for the regularised one below), i.e. about $12\%$ of the raw charge: a very large saving, but one order of magnitude, not several.
\end{coursenote}

\paragraph{Weaknesses of PCA hedging} \ts{13}{10}{48:14}\ts{13}{10}{49:19}.
\begin{itemize}
  \item The components are tuned to a historical window, and the coefficients of $R$ are not
      stable, in particular for a short window that reflects the current regime; when they change the hedge must be rebalanced, which costs money.
  \item The instability comes from the noisy components attached to
      small singular values, exactly the mechanism of \cref{ch20:sec-inverse}.
  \item PCA is a least-squares approximation and reacts strongly to outliers: a single day on which rates spiked and reverted has an unwarranted influence.
  \item It overfits history: a model that would have worked for the last three months need not work for the next three.
  \item It makes \emph{no assumption on the shape of the yield curve}: if
      the maturity labels of the swaps were scrambled, the computation would produce the same numbers, unscrambled, although 2 years is close to 1 year and lies between 1 and 5 years. Statistical structure
      without economic structure is the recurrent theme of \cref{ch05:sec-shrink} as well.
\end{itemize}

\subsection{A hedging model regularised by smoothness of the forward curve}\label{ch20:sec-reghedge}
The alternative \ts{13}{10}{50:21}: forget history and impose an \emph{economic assumption}. If nothing is known about what happens in ten years, there is no reason for the forward rate to have a spike at that date, so scenarios should be
smooth in forward-rate space \ts{13}{10}{51:25}. Let $J\in\R^{q\times n}$ be the Jacobian that translates a shift of the $n$ curve inputs into a shift of $q$ forward rates on a fine grid (for the toy above, $q=n$ and $J=M^{-1}$), and $L$ a
smoothness matrix acting on the forward grid, for instance the first-difference matrix, whose rows are $(\dots,1,-1,\dots)$. We look for the smoothest scenarios consistent with the hedging instruments: for each hedge $j$,
\begin{equation}\label{ch20:eq-regR}
  \min_{\bm\rho}\ \|LJ\bm\rho\|^2\quad\text{subject to}\quad H\T\bm\rho=\bm e_j ,
\end{equation}
where $\bm\rho=R_{\cdot j}$ is the input scenario ($1$~bp on hedge $j$, $0$ on the other hedges) and $\|LJ\bm\rho\|^2$ is the roughness of the forward curve it induces. With a soft constraint, penalising
$\|H\T\bm\rho-\bm e_j\|^2+\lambda^2\|LJ\bm\rho\|^2$ and setting the gradient to zero gives, for all columns at once,
\begin{equation}\label{ch20:eq-regRsol}
  R=\bigl(HH\T+\lambda^2\,\mathsf M\bigr)^{-1}H,\qquad \mathsf M=(LJ)\T(LJ),
\end{equation}
with translated risk $\bm x=R\T\bm r$ as before \ts{13}{10}{52:25}.

\begin{proposition}[Soft and hard constraint]\label{ch20:prop-reghedge}
Let $\mathsf M=(LJ)\T LJ$ be invertible. Then $H\T\mathsf M^{-1}H$ is invertible for $H$ of full column rank, the solution of \eqref{ch20:eq-regR} is
$R^\ast=\mathsf M^{-1}H(H\T\mathsf M^{-1}H)^{-1}$, and $(HH\T+\lambda^2\mathsf M)^{-1}H\to R^\ast$ as $\lambda\to0$. If $\mathsf M$ is singular the same limit holds, by the usual penalty-method argument, as long as $\ker\mathsf M\cap\ker H\T=\{0\}$ (as in the toy, where the kernel is the parallel shift).
\end{proposition}
\begin{proof}
Lagrange: $2\mathsf M\bm\rho=H\bm\mu$, so $\bm\rho=\mathsf M^{-1}H\bm\mu/2$ and $H\T\bm\rho=\bm e_j$ fixes $\bm\mu/2=(H\T\mathsf M^{-1}H)^{-1}\bm e_j$. For the limit use the push-through identity
$(\lambda^2\mathsf M+HH\T)^{-1}H=\mathsf M^{-1}H\,(\lambda^2I+H\T\mathsf M^{-1}H)^{-1}$ (multiply both sides by $\lambda^2\mathsf M+HH\T$ and simplify) and let $\lambda\to0$. When $\mathsf M$ is singular the matrix
$HH\T+\lambda^2\mathsf M$ is still invertible under the stated condition, and the minimiser of \eqref{ch20:eq-regR} is unique on the affine set $H\T\bm\rho=\bm e_j$; I do not prove the convergence in that case.
\end{proof}

\begin{coursenote}[Slide 23 leaves out a factor]
Slide~23 shows ``$R\sim(HH\T+\lambda^2(LJ)\T LJ)^{-1}$''. As written the matrix is $n\times n$, not $n\times k$, and does not satisfy $H\T R\approx I$: the correct expression, obtained from the two conditions on the same slide, is \eqref{ch20:eq-regRsol}, with $H$ on the right.
I read ``$\sim$'' as ``is approximately'', the exact constraint being reached for $\lambda\to0$.
\end{coursenote}

With $H\T R\approx I$ up to the size of $\lambda$ (the lecture's matrix has $1.00$ and $0.00$ on the hedging rows, the smoothness being what fills the other rows) the regularised model produces the columns of slide~24 \ts{13}{10}{52:25}. They are
\emph{hat-shaped} (peaked at the hedging tenor and decaying smoothly), whereas the PCA columns follow the historical modes; the results are in the second block of \cref{ch20:tab-hedge}.

\begin{table}[ht]
\centering\footnotesize
\begin{tabular}{lrrrrrr}
\toprule
translated risk on & 1Y & 2Y & 5Y & 10Y & 30Y & total charge\\
\midrule
PCA (slide 20) & $-426{,}757$ & 1,892,770 & $-1{,}538{,}808$ & $-268{,}764$ & 293,261 & 442,036\\
regularised (slide 24) & $-487{,}769$ & 2,082,997 & $-1{,}752{,}958$ & 78,962 & 64,483 & 446,717\\
toy of this chapter & $-331{,}027$ & 1,956,764 & $-1{,}873{,}952$ & $-92{,}944$ & 343,425 & 459,811\\
\bottomrule
\end{tabular}
\caption{Translated risk $\bm x=R\T\bm r$ of the book of \cref{ch20:tab-risk} on the five liquid swaps, and the charge at $0.1$~bp per instrument (the slides print the charge of the PCA
row and I computed the others; the toy, my computation, uses $J=M^{-1}$ of \cref{ch20:par-toy}, first differences for $L$ and $\lambda^2=10^{-4}$). Multiplying the printed two-decimal matrices by the raw risk reproduces the slide values up to the rounding of the entries (differences up to about $25{,}000$ per entry, since raw risks reach $4$ million).}\label{ch20:tab-hedge}
\end{table}

\begin{figure}[ht]
\centering
\begin{tikzpicture}
\begin{axis}[width=12cm,height=5.6cm,xlabel={maturity of the input bucket (years)},ylabel={column of $R$ for the 5Y hedge},xmin=0,xmax=31,ymin=-0.35,ymax=1.1,
  axis lines=left,grid=major,grid style={black!12},legend style={at={(0.5,-0.32)},anchor=north,legend columns=3,draw=none,font=\small},tick label style={font=\small}]
\addplot[thick,thmcol,mark=*,mark size=1.3pt] coordinates {(1,0.00) (2,0.00) (3,0.29) (4,0.70) (5,1.00) (6,0.80) (7,0.54) (8,0.33) (9,0.15) (10,0.00) (12,-0.08) (15,-0.04) (20,-0.08) (25,-0.04) (30,0.00)};
\addlegendentry{PCA (slide 20)}
\addplot[thick,defcol,mark=square*,mark size=1.3pt] coordinates {(1,0.00) (2,0.00) (3,0.44) (4,0.92) (5,1.00) (6,0.81) (7,0.58) (8,0.35) (9,0.15) (10,0.00) (12,-0.17) (15,-0.24) (20,-0.19) (25,-0.09) (30,0.00)};
\addlegendentry{regularised (slide 24)}
\addplot[thick,intcol,dashed,mark=triangle*,mark size=1.5pt] coordinates {(1,-0.00) (2,0.00) (3,0.49) (4,0.86) (5,1.00) (6,0.87) (7,0.63) (8,0.38) (9,0.16) (10,0.00) (12,-0.15) (15,-0.20) (20,-0.16) (25,-0.08) (30,-0.00)};
\addlegendentry{toy (mine)}
\end{axis}
\end{tikzpicture}
\caption{The scenario in which the 5-year hedge moves by $1$~bp and the other four hedges do not: $j=5$Y column of $R$. Both models are equal to $1$ at $5$Y and vanish at the other hedging tenors; the
regularised column is a smooth hat whose negative lobe (rates at 12--25 years move opposite to keep forwards smooth) is more pronounced than PCA's; the toy reproduces the shape. Points are the entries of the
matrices of slides 20 and 24 (two decimals).}\label{ch20:fig-hedge}
\end{figure}

\begin{finance}[Model risk in hedging and relative value]
Two hedges with different assumptions cost the same here (about $445$k) but distribute the risk differently, e.g.\ $10$Y risk of $-268{,}764$ (PCA) versus $+78{,}962$ (regularised). Traders have their own view of the dynamics of the
market, but always \emph{some} model \ts{13}{10}{47:12}. The lecture's argument for the regularised model is that it rests on a \emph{fundamental} assumption (smoothness of forwards) rather than on a window of data
\ts{13}{10}{1:26:24}. The two ideas can also be combined (a PCA regulariser inside the smoothness model).
\end{finance}

\begin{exercise}[(not from the course)]
Let $J=I_4$, $L$ the first-difference matrix ($3\times4$) and $H=(\bm e_1,\bm e_4)$ (hedges in buckets 1 and 4) in \eqref{ch20:eq-regR}. Show that the smoothest scenarios are the linear interpolations $R_{\cdot1}=(1,\tfrac23,\tfrac13,0)\T$ and $R_{\cdot2}=(0,\tfrac13,\tfrac23,1)\T$, compute the translated risk $\bm x=R\T\bm r$ for $\bm r=(0,1,1,0)\T$, and check that $\mathsf M=(LJ)\T LJ$ is singular with kernel spanned by $(1,1,1,1)\T$ (and that \cref{ch20:prop-reghedge} still applies). What would the second-difference matrix give?
\end{exercise}

\begin{exercise}[(not from the course)]
Prove the push-through identity $(\lambda^2\mathsf M+HH\T)^{-1}H=\mathsf M^{-1}H(\lambda^2I+H\T\mathsf M^{-1}H)^{-1}$ used in \cref{ch20:prop-reghedge}, and use it to show that $H\T R_\lambda\to I$ with $R_\lambda=(HH\T+\lambda^2\mathsf M)^{-1}H$ and that $H\T R_\lambda=H\T\mathsf M^{-1}H(\lambda^2I+H\T\mathsf M^{-1}H)^{-1}$ (so the deviation from $I$ is of order $\lambda^2$).
\end{exercise}

% ==================================================================
\section{Linear inverse problems, the SVD and Tikhonov regularisation}\label{ch20:sec-inverse}

The lecture analyses the calibration problem with the tools of linear algebra \ts{13}{10}{1:11:29}. This section develops that analysis in full.

\subsection{The problem and its singular value decomposition}\label{ch20:sec-svdprob}
Let $A\in\R^{M\times N}$ be the linearised pricing map (model parameters $\bm x\in\R^N$ to $M$ market observables), and
\begin{equation}\label{ch20:eq-inv}
  \bm y=A\bm x+\bm\varepsilon
\end{equation}
the market data, with an error $\bm\varepsilon$ (bid--offer, stale quotes, model error) \ts{13}{10}{1:11:29}. Take the thin SVD (\cref{thm:svd}, \cref{eq:svdsum})
$A=\sum_{i=1}^rs_i\bm u_i\bm v_i\T$ with $s_1\ge\dots\ge s_r>0$, $r=\rank A$, orthonormal $\bm u_i\in\R^M$ and $\bm v_i\in\R^N$. In the picture of the slide, $A$ is a rotation, a scaling by $s_i$ and a rotation, and the
inverse is the rotation, the scaling by $1/s_i$, the rotation \ts{13}{10}{1:11:29}.

\begin{proposition}[Minimum-norm least squares]\label{ch20:prop-pinv}
The vectors minimising $\|A\bm x-\bm y\|$ are $\bm x=\sum_{i\le r}\frac{\bm u_i\T\bm y}{s_i}\bm v_i+\bm z$ with $\bm z\in\ker A$; the one of smallest norm is
\begin{equation}\label{ch20:eq-pinv}
  \bm x^\dagger=A^\dagger\bm y=\sum_{i=1}^r\frac{\bm u_i\T\bm y}{s_i}\,\bm v_i=VS^{-1}U\T\bm y .
\end{equation}
It equals $(A\T A)^{-1}A\T\bm y$ if $r=N$ (more data than parameters), and $A\T(AA\T)^{-1}\bm y$, which fits the data exactly, if $r=M<N$.
\end{proposition}
\begin{proof}
Write $\bm x=\sum_j\xi_j\bm v_j+\bm z$ (with $\bm z\perp\bm v_j$). Then $\|A\bm x-\bm y\|^2=\sum_{i\le r}(s_i\xi_i-b_i)^2+\|\bm y\|^2-\sum_{i\le r}b_i^2$, $b_i=\bm u_i\T\bm y$, minimised for $\xi_i=b_i/s_i$; the component $\bm z\in\ker A$ does not enter and the norm is smallest for $\bm z=\bm 0$.
The two closed forms follow by substituting the SVD.
\end{proof}
The slide's formula $\bm x=(A\T A)^{-1}A\T\bm y$ is thus valid only when the columns of $A$ are independent; for the surface calibration, where there are fewer instruments than parameters,
one needs \eqref{ch20:eq-pinv} (the SVD form on the same slide is general).

\subsection{Noise amplification and the condition number}\label{ch20:sec-noise}
Insert \eqref{ch20:eq-inv} into \eqref{ch20:eq-pinv}: $\bm x^\dagger=A^\dagger A\bm x+\sum_i\frac{\bm u_i\T\bm\varepsilon}{s_i}\bm v_i$, the second term being the \emph{propagated error}
\ts{13}{10}{1:12:32}:
\[
  \|\delta\bm x\|^2=\sum_{i=1}^r\frac{(\bm u_i\T\bm\varepsilon)^2}{s_i^2},\qquad
  \E[\|\delta\bm x\|^2]=\sigma^2\sum_{i=1}^r\frac1{s_i^2}\quad\text{if }\Cov\bm\varepsilon=\sigma^2I .
\]
Each mode $i$ of the noise is divided by $s_i$: a small singular value turns an innocent perturbation of the data into a large variation of the parameters. A scale-free measure is the
\emph{condition number} $\kappa=s_1/s_r$ \ts{13}{10}{1:12:32}: for $\bm x$ in the row space of $A$,
\[
  \frac{\|\delta\bm x\|}{\|\bm x\|}\le\kappa\,\frac{\|\bm\varepsilon\|}{\|A\bm x\|},
\]
because $\|\delta\bm x\|\le\|\bm\varepsilon\|/s_r$ and $\|A\bm x\|\le s_1\|\bm x\|$; the bound is attained for $\bm\varepsilon\parallel\bm u_r$, $\bm x\parallel\bm v_1$. A problem is \emph{well posed}
(Hadamard) if a solution exists, is unique and depends continuously on the data; it is \emph{ill-posed} if one of the three fails, and \emph{ill-conditioned} if the discretised problem has a large $\kappa$. The
lecture's calibration problems are underdetermined (uniqueness fails: fewer instruments than parameters, hence $\ker A\ne\{0\}$) \emph{and} unstable (small $s_i$)
\ts{13}{10}{59:53}\ts{13}{10}{1:00:55}.

\begin{intuition}[A physicist's reading: the deconvolution analogue]
The typical forward map averages (the price of a swaption is an integral of the volatility surface over a region), and averaging damps the fine details of the parameters: their contribution to $\bm y$ is $s_i\xi_i$ with $s_i$ small. Inverting means amplifying
back what the data barely contain; when the data have a noise floor $\sigma$, all modes with $s_i|\xi_i|\lesssim\sigma$ are lost, and dividing by their $s_i$ turns them into noise.
This is the same as the inversion of a heat-kernel smoothing, or an image deconvolution, and the standard fix is the same.
\end{intuition}

\paragraph{The discrete Picard picture.} For noiseless data $\bm y=A\bm x$ one has $\bm u_i\T\bm y=s_i\,\bm v_i\T\bm x$: the coefficients $|\bm u_i\T\bm y|$ decay at least as fast as $s_i$ if $\bm x$ is reasonable. With noise, the coefficients level off at the noise floor
$\sigma$ and the ratios $|\bm u_i\T\bm y|/s_i$ that enter \eqref{ch20:eq-pinv} grow. \Cref{ch20:fig-spectral}(b) shows this for the toy calibration of \cref{ch20:sec-lambda}.

\begin{coursenote}[The ``noiseless'' slide]\label{ch20:cn-noiseless}
Slide~46 argues that in the noiseless case ``$A(A\T A)^{-1}A\T\bm y=\bm y$'' so that ``the model appears accurate'', but ``pricing of the actual portfolio with the model $B$ may be unstable''. The point is right; the printed identity is not: slide~46
writes ``$A\T(A\T A)^{-1}A=I$'', which is not even conformable. The correct statement is that $AA^\dagger=\sum_{i\le r}\bm u_i\bm u_i\T$ is the orthogonal projector onto the range of $A$, which is the identity on $\R^M$ exactly when $r=M$, i.e. when
the instruments are independent, for any $N\ge M$. Then \emph{every} data vector, noise included, is repriced exactly: exact repricing certifies nothing, because it is achieved by fitting the noise. What matters is the price $\bm p=B\bm x$ of the
instruments \emph{not} used in the calibration, and $B\bm x^\dagger$ inherits the amplified error $B\,\delta\bm x$. Slide~47 has the same misprint for the regularised repricing matrix, see \eqref{ch20:eq-influence}.
\end{coursenote}

\begin{exercise}[(not from the course)]
Consider $A=\begin{psmallmatrix}1&1\\1&1.01\end{psmallmatrix}$.
\begin{enumerate}[(a)]
  \item Solve $A\bm x=\bm y$ for $\bm y=(2,2.01)\T$ and for $\bm y'=(2,2.02)\T$; by what fraction do data and solution change?
  \item Compute the singular values and $\kappa(A)$ (about $4\times10^2$) and
      check that the change of (a) is compatible with the bound of \cref{ch20:sec-noise}.
  \item Compute the Tikhonov solutions \eqref{ch20:eq-tikh} for $\lambda=0.1$ for both $\bm y$ and $\bm y'$ and compare.
\end{enumerate}
\end{exercise}

\subsection{Truncated SVD and Tikhonov regularisation}\label{ch20:sec-tikh}
Both methods multiply each term of \eqref{ch20:eq-pinv} by a \emph{filter factor} $w_i\in[0,1]$:
\begin{equation}\label{ch20:eq-filter}
  \bm x_w=\sum_{i=1}^rw_i\,\frac{\bm u_i\T\bm y}{s_i}\,\bm v_i .
\end{equation}
\begin{itemize}
  \item \emph{Truncated SVD} (TSVD): $w_i=1$ for $i\le k$ and $0$ for $i>k$ \ts{13}{10}{1:16:49}. One keeps the well-determined modes and discards the rest, which are the part of
        the parameter space that the calibration cannot see (the ``null space'' of the model, in the speaker's words). It is the same idea as the PCA risk model of \cref{ch20:sec-pcahedge} and is principal-components regression in the language of \cref{ch05:sec-shrink}.
        The danger: if the discarded modes do affect the price of the portfolio, the truncation is a bias one pays for.
  \item \emph{Tikhonov regularisation} \ts{13}{10}{1:14:44}:
        \begin{equation}\label{ch20:eq-tikh}
          \bm x_\lambda=\arg\min_{\bm x}\Bigl\{\|A\bm x-\bm y\|^2+\lambda^2\|\bm x\|^2\Bigr\}=(A\T A+\lambda^2I)^{-1}A\T\bm y=\sum_{i=1}^r\underbrace{\frac{s_i^2}{s_i^2+\lambda^2}}_{w_i}\frac{\bm u_i\T\bm y}{s_i}\bm v_i .
        \end{equation}
\end{itemize}
The gradient of the objective is $2A\T(A\bm x-\bm y)+2\lambda^2\bm x$; setting it to zero gives the normal equations $(A\T A+\lambda^2I)\bm x=A\T\bm y$, uniquely solvable for every $\lambda>0$ even if $A$ is rank-deficient (the matrix is
positive definite), and $A\T=V S U\T$ gives the filter form, since $(A\T A+\lambda^2I)^{-1}A\T=V\,\diag\!\bigl(\tfrac{s_i}{s_i^2+\lambda^2}\bigr)U\T$. Hence Tikhonov is a \emph{smooth} filter: $w_i\approx1$ for $s_i\gg\lambda$ and $w_i\approx s_i^2/\lambda^2$ for $s_i\ll\lambda$
(\cref{ch20:fig-spectral}(a)). The essential point is stated by the speaker: one no longer divides by a small number \ts{13}{10}{1:14:44}. Quantitatively, the gain applied to a mode of noise is
$\frac{w_i}{s_i}=\frac{s_i}{s_i^2+\lambda^2}\le\frac1{2\lambda}$ (maximum at $s_i=\lambda$), instead of $1/s_r$ for the exact inverse, and $1/s_k$ for TSVD.

\begin{intuition}[Ridge regression in another language]
This is the ridge regression of \cref{ch05:sec-shrink} with $d_j\to s_i$ and ridge parameter $\lambda_{\rm ridge}=\lambda^2$ (the lecture writes the penalty as $\|\lambda\bm x\|^2$). The smooth filter $s^2/(s^2+\lambda^2)$ is the same shrinkage factor,
and \cref{ch05:prop-ridgebayes} gives its Bayesian reading: a Gaussian prior $\bm x\sim N(\bm 0,\tau^2I)$ with $\lambda^2=\sigma^2/\tau^2$. Early stopping of gradient descent acts as another filter of the same family (\cref{ch16:sec-gd}).
\end{intuition}

The price of stability is that the calibration is no longer exact. The fit and the ``re-pricing'' operator are, by \eqref{ch20:eq-tikh},
\begin{equation}\label{ch20:eq-influence}
  A\bm x_\lambda=H_\lambda\bm y,\qquad H_\lambda=A(A\T A+\lambda^2I)^{-1}A\T=\sum_iw_i\,\bm u_i\bm u_i\T\ne I,
\end{equation}
the corrected form of the ``$A\T(A\T A+\lambda^2I)^{-1}A\neq I$'' of slide~47 \ts{13}{10}{1:15:49}; the residual is $(I-H_\lambda)\bm y$. Its trace $\operatorname{tr}H_\lambda=\sum_iw_i$ is the \emph{effective number of parameters} used by the fit.

\begin{figure}[ht]
\centering
\begin{tikzpicture}
\begin{axis}[name=a,width=6.2cm,height=5cm,xlabel={$\log_{10}(s/\lambda)$},ylabel={filter $w$},xmin=-2,xmax=2,ymin=-0.05,ymax=1.1,
  axis lines=left,grid=major,grid style={black!12},tick label style={font=\small},legend style={at={(0.5,-0.32)},anchor=north,draw=none,legend columns=2,font=\small}]
\addplot[thick,thmcol,domain=-2:2,samples=80] {10^(2*x)/(10^(2*x)+1)};
\addlegendentry{Tikhonov}
\addplot[thick,defcol,dashed] coordinates {(-2,0)(0,0)(0,1)(2,1)};
\addlegendentry{TSVD}
\node[anchor=north west] at (axis cs:-1.95,1.08) {\small(a)};
\end{axis}
\begin{axis}[at={(a.east)},anchor=west,xshift=1.7cm,width=6.2cm,height=5cm,xlabel={mode $i$},ymode=log,ymin=5e-4,ymax=1,xmin=0.5,xmax=10.5,
  axis lines=left,grid=major,grid style={black!12},tick label style={font=\small},legend style={at={(0.5,-0.32)},anchor=north,draw=none,legend columns=3,font=\small}]
\addplot[thick,thmcol,mark=*,mark size=1.3pt] coordinates {(1,0.35564) (2,0.30201) (3,0.27821) (4,0.21286) (5,0.15147) (6,0.10928) (7,0.07985) (8,0.04535) (9,0.02417) (10,0.00281)};
\addlegendentry{$s_i$}
\addplot[thick,defcol,mark=square*,mark size=1.3pt] coordinates {(1,0.45418) (2,0.49752) (3,0.12507) (4,0.00788) (5,0.00641) (6,0.00865) (7,0.00814) (8,0.00214) (9,0.00357) (10,0.00174)};
\addlegendentry{$|\bm u_i\T\bm y|$}
\addplot[thick,intcol,dashed] coordinates {(0.5,0.00189)(10.5,0.00189)};
\addlegendentry{noise}
\node[anchor=south west] at (axis cs:0.8,6e-4) {\small(b)};
\end{axis}
\end{tikzpicture}
\caption{(a) Filter factors of Tikhonov ($w=s^2/(s^2+\lambda^2)$) and of TSVD (cut at $s=\lambda$). (b) Singular values $s_i$ of the toy calibration matrix of \cref{ch20:sec-lambda} (ten instruments, 100 parameters, $\kappa=126$) and the coefficients $|\bm u_i\T\bm y|$ of noisy data
(1\% price noise; my computation). From mode 4 on the coefficients stop tracking $s_i$ and stay within one to four times the noise floor $\|\bm\varepsilon\|/\sqrt M$ (dashed), whereas $s_i$ keeps falling: the terms $\bm u_i\T\bm y/s_i$ of \eqref{ch20:eq-pinv} are then dominated by noise.}\label{ch20:fig-spectral}
\end{figure}

\begin{exercise}[(not from the course)]
Prove the following for Tikhonov regularisation with $L=I$:
\begin{enumerate}[(a)]
  \item $\rho(\lambda)=\|A\bm x_\lambda-\bm y\|$ is non-decreasing and $\|\bm x_\lambda\|$ non-increasing in $\lambda$;
  \item as $\lambda\to0$, $\bm x_\lambda\to A^\dagger\bm y$ (\cref{ch20:prop-pinv}), also when $A$ is rank-deficient;
  \item as $\lambda\to\infty$, $\bm x_\lambda\approx A\T\bm y/\lambda^2$;
  \item $0\le\operatorname{tr}H_\lambda\le r$ and it decreases from $r$ to $0$.
\end{enumerate}
\end{exercise}

\subsection{General form: weights, smoothness matrices and several penalties}\label{ch20:sec-general}
Penalising the amplitude $\|\bm x\|$ is rarely what one wants. The general problem, with a diagonal weight matrix $W$ and a penalty matrix $L$, is
\begin{equation}\label{ch20:eq-gen}
  \bm x_\lambda=\arg\min_{\bm x}\Bigl\{\|W(A\bm x-\bm y)\|^2+\lambda^2\|L\bm x\|^2\Bigr\}=(A\T W^2A+\lambda^2L\T L)^{-1}A\T W^2\bm y .
\end{equation}
The solution is unique iff $\ker(WA)\cap\ker L=\{0\}$. It is an ordinary least-squares problem for the \emph{augmented system}
$\begin{psmallmatrix}WA\\ \lambda L\end{psmallmatrix}\bm x\approx\begin{psmallmatrix}W\bm y\\ \bm 0\end{psmallmatrix}$, which is how it is solved stably in practice (QR, \cref{ch05:sec-qr}). Three remarks.
\begin{itemize}
  \item \emph{Weights.} $W=\diag(1/\sigma_i)$, with $\sigma_i$ the tolerance of instrument $i$ (its bid--offer): the data term is then a $\chi^2$ and a residual of one tolerance costs one unit. Instruments that are
        important or liquid get a small tolerance and are fitted more tightly \ts{13}{10}{1:07:09}.
  \item \emph{Smoothness instead of amplitude.} With $L=I$ the penalty pulls $\bm x$ towards $\bm 0$, which for a volatility surface means a volatility of zero. The speaker's point \ts{13}{10}{1:19:00}: ask first what one does \emph{not} like. A perfectly flat
        surface is acceptable, so the penalty should vanish on it, i.e. $L$ should have the constants in its kernel: the first-difference matrix with rows $(\dots,1,-1,\dots)$ of slide~41 penalises $\sum_i(x_i-x_{i+1})^2$, a discrete $\int|\nabla x|^2$ (in two dimensions the differences of neighbouring nodes along both axes).
  \item \emph{No bias on the kernel of $L$.} If $A\bm x=\bm y$ exactly and $L\bm x=\bm 0$, then $\bm x_\lambda=\bm x$: indeed $\bm x_\lambda=(A\T W^2A+\lambda^2L\T L)^{-1}A\T W^2A\bm x=\bm x-\lambda^2(A\T W^2A+\lambda^2L\T L)^{-1}L\T L\bm x=\bm x$.
        Regularisation biases the solution only in the directions it penalises, so the bias is small when the penalty encodes true prior knowledge \ts{13}{10}{1:17:52}\ts{13}{10}{1:19:00}.
\end{itemize}
Slide~42 combines two penalties, one on the \emph{change} of the surface with respect to a prior $\bm v_0$ (matrix $L_1$ acting on $\bm v-\bm v_0=B\bm x$) and one of Tikhonov type on the coefficients (matrix $L_2$), with parameters $\lambda_1,\lambda_2$ (\cref{ch20:sec-volsolve}); this is \eqref{ch20:eq-gen} with
$\lambda^2L\T L\to\lambda_1^2(L_1B)\T L_1B+\lambda_2^2L_2\T L_2$. The Bayesian reading is the same as for ridge: a Gaussian prior with precision $(\lambda^2/\sigma^2)L\T L$, a random-walk prior for a difference penalty.

\begin{coursenote}[Notation of the slides]
On slide~42 the inverse matrix is called $A$, which is also the model matrix of slides 44--48 (here $K=(J\T W^2J+\dots)^{-1}$ in \cref{ch20:sec-volsolve}); and the penalty matrix $L_2$ is applied to the volatility vector $\bm v$ on slide~40 but to the coefficients $\bm x$ on slide~42.
\end{coursenote}

\subsection{Bias, variance and the meaning of $\lambda$}\label{ch20:sec-biasvar}
For Tikhonov with $L=I$ the error can be computed exactly. Let $\bm y=A\bm x+\bm\varepsilon$, $\E[\bm\varepsilon]=\bm 0$, $\Cov\bm\varepsilon=\sigma^2I$, and $\xi_i=\bm v_i\T\bm x$. From \eqref{ch20:eq-filter}, $\bm x_\lambda=\sum_iw_i(\xi_i+\bm u_i\T\bm\varepsilon/s_i)\bm v_i$, hence
\begin{equation}\label{ch20:eq-mse}
  \E[\|\bm x_\lambda-\bm x\|^2]=\sum_{i=1}^r\Bigl[\underbrace{(1-w_i)^2\xi_i^2}_{\text{bias}^2}+\underbrace{\sigma^2\frac{w_i^2}{s_i^2}}_{\text{variance}}\Bigr]+\sum_{i>r}\xi_i^2 ,\qquad 1-w_i=\frac{\lambda^2}{s_i^2+\lambda^2},\quad\frac{w_i^2}{s_i^2}=\frac{s_i^2}{(s_i^2+\lambda^2)^2}.
\end{equation}
(I checked \eqref{ch20:eq-mse} by Monte Carlo on a random $8\times6$ matrix: $0.7747$ against $0.7750$ from $4\times10^4$ samples.) The last sum is the part of $\bm x$ in the kernel of $A$, which no method can recover.
Bias grows and variance falls with $\lambda$; at $\lambda=0$ the variance $\sigma^2\sum s_i^{-2}$ is the amplification of \cref{ch20:sec-noise}. Minimising each term over $w_i$: $-2(1-w)\xi^2+2w\sigma^2/s^2=0$ gives
\begin{equation}\label{ch20:eq-wiener}
  w_i^\ast=\frac{s_i^2\xi_i^2}{s_i^2\xi_i^2+\sigma^2}=\frac{s_i^2}{s_i^2+\sigma^2/\xi_i^2},
\end{equation}
a Wiener filter, i.e.\ Tikhonov with a mode-dependent $\lambda_i^2=\sigma^2/\xi_i^2$: a mode is kept if the signal it contributes to the data, $s_i|\xi_i|$, exceeds the noise $\sigma$, and suppressed otherwise. TSVD is the $0/1$ version of this filter and a single $\lambda^2=\sigma^2/\tau^2$
(with $\tau^2$ a typical size of $\xi_i^2$) is the ridge--Bayes value of \cref{ch05:prop-ridgebayes}. The optimum cannot be used as it stands, since $\xi_i$ is unknown, but it says what $\lambda$ \emph{is}: a noise-to-signal ratio for the parameters.

\begin{exercise}[(not from the course)]
Take $N=M=1$: $y=ax+\varepsilon$, $\varepsilon\sim N(0,\sigma^2)$. Find the Tikhonov estimator, its bias and variance, and the $\lambda$ that minimises the mean-square error as a function of $x$; compare with \eqref{ch20:eq-wiener}. What happens if one substitutes the ridge--Bayes value $\lambda^2=\sigma^2/\tau^2$ and the true $|x|$ is much larger or much smaller than $\tau$?
\end{exercise}

% ==================================================================
\section{Choosing the regularisation parameter (not covered in class)}\label{ch20:sec-lambda}
The lecture calls $\lambda$ a ``small regularisation parameter'' and says nothing on how to fix it \ts{13}{10}{52:25}. Three standard criteria, then a numerical test. Let $\rho(\lambda)=\|W(A\bm x_\lambda-\bm y)\|$ be the residual.
For $L=I$ and $b_i=\bm u_i\T\bm y$, $\rho^2=\sum_{i\le r}(1-w_i)^2b_i^2+\text{const}$ increases with $\lambda$ and $\|\bm x_\lambda\|^2=\sum_iw_i^2b_i^2/s_i^2$ decreases (each $w_i$ decreases with $\lambda$), so the trade-off is monotone.

\begin{itemize}
  \item \emph{Discrepancy principle} (Morozov): choose $\lambda$ with $\rho(\lambda)=\delta$, the size of the data error, $\delta\approx\sigma\sqrt M$ (with $W=\diag(1/\sigma_i)$: $\rho^2\approx M$). Rationale: fitting closer than the noise fits the noise. This is the lecture's remark that every input has a
        tolerance and calibrating exactly is pointless \ts{13}{10}{1:06:08}; it needs the noise level, here the bid--offer.
  \item \emph{Generalised cross-validation.} Leave-one-out cross-validation predicts each datum from the others. For a linear smoother $\hat{\bm y}=H_\lambda\bm y$ one has $y_i-\hat y_i^{(-i)}=(y_i-\hat y_i)/(1-(H_\lambda)_{ii})$ (\emph{proof}: the fit computed without datum $i$ is also the minimiser when
        $y_i$ is replaced by its own prediction $\hat y_i^{(-i)}$, since that term is then zero; so $\hat y_i^{(-i)}=\sum_{j\ne i}H_{ij}y_j+H_{ii}\hat y_i^{(-i)}$). Replacing $(H_\lambda)_{ii}$ by the average $\operatorname{tr}H_\lambda/M$ gives
        \[
          \mathrm{GCV}(\lambda)=\frac{M\,\rho(\lambda)^2}{\bigl(M-\operatorname{tr}H_\lambda\bigr)^2},
        \]
        which needs no noise level; it is the counterpart of the cross-validation used for ridge in \cref{ch05:sec-shrink}.
  \item \emph{L-curve.} Plot $\log\rho(\lambda)$ against $\log\|L\bm x_\lambda\|$; the corner of the ``L'' separates the regime dominated by noise (small $\lambda$, large seminorm) from the regime dominated by bias.
\end{itemize}

\begin{casestudy}[A toy calibration with regularisation (my computation; not in class)]\label{ch20:cs-toy}
\textbf{Goal.} Mimic the calibration of a volatility slice. \textbf{Set-up.} A smooth ``true'' volatility $v(s)$ on a grid of $N=100$ points, $s\in[0,10]$ (level $0.19$--$0.25$); $M=10$ instruments, each an average of $v$ with a Gaussian window (centres $0.5$--$9.7$, widths growing from $0.46$ to $1.56$, overlapping); prices $y_i=(A\bm v)_i(1+0.01z_i)$ with $z_i$ standard normal, so $\delta=\|\bm\varepsilon\|=0.0060$;
$\kappa(A)=126$ on the ten modes.
\textbf{Methods.}
\begin{enumerate}[(a)]
  \item The lecture's ``formal'' solution: ten piecewise-constant basis functions, square Jacobian, exact inversion.
  \item Tikhonov \eqref{ch20:eq-gen} with the first-difference matrix $L$ of slide~41, $W=I$, on the fine grid.
\end{enumerate}
\textbf{Results} (\cref{ch20:fig-toy}).
\begin{enumerate}[(a)]
  \item Reprices the ten inputs to $5\times10^{-17}$ but the surface swings between $0.002$ and $0.40$ (root-mean-square error $0.106$, roughness $\|L\bm v\|=0.64$), and a $1\%$ change in one input moves it by $16\%$ of its level in the worst place: the skyline of the
      lecture \ts{13}{10}{1:04:01}.
  \item With $\lambda$ from GCV ($0.44$, effective number of parameters $\operatorname{tr}H_\lambda=6.2$): error $0.0031$, roughness $0.0144$ (truth $0.0147$), the same $1\%$ perturbation moves it by $1.3\%$ of its level. The discrepancy principle
      gives $\lambda=0.96$ and error $0.0032$; the best possible value (knowing the truth) is $0.0031$ at $\lambda=0.65$. GCV picks a residual of $0.46\,\delta$, a mild over-fit, typical of GCV.
\end{enumerate}
\textbf{Which penalty?} With the same tuning, the best error is $0.0167$ for $L=I$ (amplitude), $0.0039$ for the second difference and $0.0031$ for the first difference. TSVD of $A$ (minimum-norm solution) reaches at best $0.018$ (with $k=8$ modes) and
$0.063$ with all ten: an unsuitable penalty (here: the minimum norm, which pulls the surface towards zero) gives errors five to six times larger than the suitable one, whereas $\lambda$ can be changed by a factor two with hardly any effect on the error.
\textbf{Lesson.} Regularisation is a matter of choosing \emph{what} to penalise (prior knowledge of the solution) first, and how much second.
\end{casestudy}

\begin{figure}[ht]
\centering
\begin{tikzpicture}
\begin{axis}[name=a,width=6.4cm,height=5.4cm,xlabel={forward time $s$},ylabel={volatility},xmin=0,xmax=10,ymin=-0.02,ymax=0.42,
  axis lines=left,grid=major,grid style={black!12},tick label style={font=\small},legend style={at={(0.5,-0.3)},anchor=north,draw=none,legend columns=1,font=\small}]
\addplot[thick,black,dashed] coordinates {(0.000,0.2309) (0.303,0.2305) (0.606,0.2305) (0.909,0.2312) (1.212,0.2333) (1.515,0.2368) (1.818,0.2416) (2.121,0.2467) (2.424,0.2509) (2.727,0.2528) (3.030,0.2511) (3.333,0.2457) (3.636,0.2371) (3.939,0.2267) (4.242,0.2159) (4.545,0.2063) (4.848,0.1986) (5.152,0.1931) (5.455,0.1895) (5.758,0.1876) (6.061,0.1868) (6.364,0.1868) (6.667,0.1873) (6.970,0.1882) (7.273,0.1894) (7.576,0.1907) (7.879,0.1922) (8.182,0.1937) (8.485,0.1953) (8.788,0.1969) (9.091,0.1984) (9.394,0.1999) (9.697,0.2012) (10.000,0.2024)};
\addlegendentry{truth}
\addplot[thick,thmcol,const plot] coordinates {(0.000,0.2275) (0.101,0.2275) (0.202,0.2275) (0.303,0.2275) (0.404,0.2275) (0.505,0.2275) (0.606,0.2275) (0.707,0.2275) (0.808,0.2275) (0.909,0.2275) (1.010,0.2431) (1.111,0.2431) (1.212,0.2431) (1.313,0.2431) (1.414,0.2431) (1.515,0.2431) (1.616,0.2431) (1.717,0.2431) (1.818,0.2431) (1.919,0.2431) (2.020,0.2528) (2.121,0.2528) (2.222,0.2528) (2.323,0.2528) (2.424,0.2528) (2.525,0.2528) (2.626,0.2528) (2.727,0.2528) (2.828,0.2528) (2.929,0.2528) (3.030,0.2457) (3.131,0.2457) (3.232,0.2457) (3.333,0.2457) (3.434,0.2457) (3.535,0.2457) (3.636,0.2457) (3.737,0.2457) (3.838,0.2457) (3.939,0.2457) (4.040,0.1940) (4.141,0.1940) (4.242,0.1940) (4.343,0.1940) (4.444,0.1940) (4.545,0.1940) (4.646,0.1940) (4.747,0.1940) (4.848,0.1940) (4.949,0.1940) (5.051,0.2413) (5.152,0.2413) (5.253,0.2413) (5.354,0.2413) (5.455,0.2413) (5.556,0.2413) (5.657,0.2413) (5.758,0.2413) (5.859,0.2413) (5.960,0.2413) (6.061,0.0444) (6.162,0.0444) (6.263,0.0444) (6.364,0.0444) (6.465,0.0444) (6.566,0.0444) (6.667,0.0444) (6.768,0.0444) (6.869,0.0444) (6.970,0.0444) (7.071,0.3991) (7.172,0.3991) (7.273,0.3991) (7.374,0.3991) (7.475,0.3991) (7.576,0.3991) (7.677,0.3991) (7.778,0.3991) (7.879,0.3991) (7.980,0.3991) (8.081,0.0020) (8.182,0.0020) (8.283,0.0020) (8.384,0.0020) (8.485,0.0020) (8.586,0.0020) (8.687,0.0020) (8.788,0.0020) (8.889,0.0020) (8.990,0.0020) (9.091,0.2893) (9.192,0.2893) (9.293,0.2893) (9.394,0.2893) (9.495,0.2893) (9.596,0.2893) (9.697,0.2893) (9.798,0.2893) (9.899,0.2893) (10.000,0.2893)};
\addlegendentry{formal (exact repricing)}
\addplot[thick,defcol] coordinates {(0.000,0.2265) (0.303,0.2273) (0.606,0.2297) (0.909,0.2338) (1.212,0.2387) (1.515,0.2436) (1.818,0.2478) (2.121,0.2507) (2.424,0.2520) (2.727,0.2512) (3.030,0.2483) (3.333,0.2431) (3.636,0.2359) (3.939,0.2274) (4.242,0.2183) (4.545,0.2093) (4.848,0.2013) (5.152,0.1946) (5.455,0.1896) (5.758,0.1862) (6.061,0.1842) (6.364,0.1834) (6.667,0.1835) (6.970,0.1842) (7.273,0.1854) (7.576,0.1870) (7.879,0.1891) (8.182,0.1916) (8.485,0.1944) (8.788,0.1973) (9.091,0.2001) (9.394,0.2025) (9.697,0.2043) (10.000,0.2050)};
\addlegendentry{Tikhonov, $\lambda=0.44$}
\node[anchor=north] at (axis cs:8,0.4) {\small(a)};
\end{axis}
\begin{axis}[at={(a.east)},anchor=west,xshift=1.7cm,width=6.4cm,height=5.4cm,xlabel={$\log_{10}\lambda$},xmin=-2,xmax=1,ymode=log,ymin=0.1,ymax=100,
  axis lines=left,grid=major,grid style={black!12},tick label style={font=\small},legend style={at={(0.5,-0.3)},anchor=north,draw=none,legend columns=1,font=\small}]
\addplot[thick,thmcol] coordinates {(-4.000,20.0356) (-3.831,20.0356) (-3.661,20.0324) (-3.492,20.0227) (-3.322,20.0000) (-3.153,19.9547) (-2.983,19.8576) (-2.814,19.6505) (-2.644,19.2071) (-2.475,18.3139) (-2.305,16.6343) (-2.136,13.8900) (-1.966,10.3042) (-1.797,6.8123) (-1.627,4.3398) (-1.458,3.0388) (-1.288,2.4207) (-1.119,2.0097) (-0.949,1.6214) (-0.780,1.2945) (-0.610,1.1003) (-0.441,1.0259) (-0.271,1.0000) (-0.102,1.0032) (0.068,1.1586) (0.237,1.6570) (0.407,2.4207) (0.576,3.1812) (0.746,3.8576) (0.915,4.5955)};
\addlegendentry{error / best error}
\addplot[thick,defcol,dashed] coordinates {(-4.000,5.1023) (-3.831,5.1023) (-3.661,5.1023) (-3.492,5.1023) (-3.322,5.1023) (-3.153,5.1023) (-2.983,5.1023) (-2.814,5.1004) (-2.644,5.0967) (-2.475,5.0911) (-2.305,5.0763) (-2.136,5.0428) (-1.966,4.9740) (-1.797,4.8326) (-1.627,4.5536) (-1.458,4.0644) (-1.288,3.3538) (-1.119,2.5800) (-0.949,1.9438) (-0.780,1.4953) (-0.610,1.1830) (-0.441,1.0000) (-0.271,1.0327) (-0.102,1.5727) (0.068,3.3036) (0.237,7.1168) (0.407,12.9092) (0.576,19.3080) (0.746,26.6555) (0.915,38.3557)};
\addlegendentry{GCV / minimum}
\addplot[thick,intcol] coordinates {(-4.000,0.0000) (-3.831,0.0167) (-3.661,0.0167) (-3.492,0.0500) (-3.322,0.0833) (-3.153,0.2000) (-2.983,0.4333) (-2.814,0.9333) (-2.644,1.9833) (-2.475,4.1333) (-2.305,8.1500) (-2.136,14.7500) (-1.966,23.4500) (-1.797,32.1333) (-1.627,38.7167) (-1.458,42.8667) (-1.288,45.5333) (-1.119,48.2167) (-0.949,52.2000) (-0.780,57.6333) (-0.610,63.7833) (-0.441,71.3500) (-0.271,85.9667) (-0.102,122.8000) (0.068,201.8500) (0.237,330.2000) (0.407,487.9833) (0.576,645.5000) (0.746,809.7833) (0.915,1024.7500)};
\addlegendentry{residual $\rho/\delta$}
\addplot[gray,dotted] coordinates {(-2,1)(1,1)};
\node[anchor=north west] at (axis cs:-1.95,90) {\small(b)};
\end{axis}
\end{tikzpicture}
\caption{The toy calibration of \cref{ch20:cs-toy} (my computation). (a) The formal exact solution is the ``skyline''; the smoothness-regularised one follows the truth to $0.003$. (b) Error, GCV and residual against $\lambda$ (each normalised as indicated; dotted line: $\rho=\delta$, the discrepancy principle). The minima of the error and of GCV are broad and lie within a factor two of each other.}\label{ch20:fig-toy}
\end{figure}

\begin{exercise}[(not from the course)]
Reproduce the toy of \cref{ch20:cs-toy} (any language): choose your own smooth $v$ on $[0,10]$, ten Gaussian-window instruments, $1\%$ noise; compute the formal solution, the Tikhonov solution for a grid of $\lambda$, the discrepancy and GCV parameters. Vary the penalty ($L=I$, first difference, second difference) and the noise level, and report which choice is most sensitive to $\lambda$.
\end{exercise}

% ==================================================================
\section{Calibrating a volatility surface}\label{ch20:sec-vol}

\subsection{The pricing engine and the HJM model}\label{ch20:sec-hjmvol}
The lecture now steps back \ts{13}{10}{53:26}. A \emph{pricing engine} takes the trader's portfolio and the model parameters (curves, volatility surface, \dots) and returns values and risks. To make
the parameters consistent with the market, a \emph{calibration engine} compares the model prices of a set of benchmark instruments with their market prices and adjusts the parameters until they agree
(slide~25); only then is the portfolio priced. The example is the Heath--Jarrow--Morton model (\cref{ch:hjm}), used to price volatility products by Monte Carlo \ts{13}{10}{54:35}:
\begin{equation}\label{ch20:eq-hjm}
  \dd f_{t,s}=\mu_{t,s}\dd t+f_{t,s}^{\beta}\,V(t,s)\,\rho(t,s)\cdot\dd B^{\Q}_t ,
\end{equation}
for the forward rate $f_{t,s}$ observed at calendar time $t$ for forward time $s$: $\mu$ is the drift (fixed by absence of arbitrage, \cref{ch:hjm}), $\beta$ the ``skew factor'' ($\beta=1$ lognormal, $\beta=0$ normal model), $\rho$ the correlation and factor structure, $B^\Q$ a
Brownian motion, and $V(t,s)$ the \emph{volatility surface} that is the subject of the lecture \ts{13}{10}{55:41}. The Monte Carlo advances the whole forward curve from one quarter to the next, and every step needs the value of $V$ on
the corresponding cell of the triangular grid of (calendar time, forward time) \ts{13}{10}{56:47}; the instruments that pin it down, caps and swaptions, are sensitive to \emph{regions} of this triangle (\ts{13}{10}{57:48}; a swaption
depends on the forward rates that will be observed at the option expiry and on their volatility until then), and the regions of different instruments overlap heavily \ts{13}{10}{58:50}.

\subsection{Why the problem is hard}\label{ch20:sec-volhard}
The surface has a value for each of the cells of a triangle with 240 quarters (60 years) on a side, that is $240\cdot241/2=28{,}920$ unknowns \ts{13}{10}{59:53} (the slide says ``$\sim$28k''), whereas only $20$--$50$ swaptions and caps are quoted. The problem is underdetermined
(the kernel of the forward map has dimension at least $28{,}870$), ill-posed because of the overlapping sensitivities, and its unregularised solution is unstable and noisy \ts{13}{10}{1:00:55}. A full $28{,}920\times28{,}920$ matrix of doubles takes $6.7$~GB (my count),
which the speaker rules out as ``no memory''. And the result must be \emph{smooth}: two similar swaptions in the portfolio should get similar prices, and a spike of volatility in some region of the future is either explained by an economic event or is an artefact
\ts{13}{10}{1:00:55}.

\begin{coursenote}[Counting the elements of the triangle]
In the transcript \ts{13}{10}{59:53} the number of triangle elements is first said to be ``200K'' and then a matrix of $28$K$\times28$K is discussed; slide~29 gives ``$\sim$28k'', which is the right count: half of $240^2=57{,}600$ plus the diagonal is $28{,}920$.
\end{coursenote}

\subsection{The ``formal'' solution and why it fails}\label{ch20:sec-volformal}
Represent the surface, flattened to a vector $\bm v$ with one entry per cell, as a prior value plus a combination of basis functions \ts{13}{10}{1:01:56}:
\begin{equation}\label{ch20:eq-vbasis}
  \bm v=\bm v_0+B\bm x,
\end{equation}
where the columns of $B$ are the basis functions (values on the grid) and $\bm x$ the weights. The pricing model gives the sensitivities of the prices of the calibration instruments to a perturbation of the surface along each basis function \ts{13}{10}{1:02:59}; working with
log-prices, $q_i=\ln(q_i^{\rm mdl}/q_i^{0})$, $q_i^{\rm in}=\ln(q_i^{\rm mkt}/q_i^{0})$,
\begin{equation}\label{ch20:eq-jac}
  J_{ij}=\frac{\partial q_i}{\partial x_j},\qquad \bm q\approx J\bm x .
\end{equation}
The \emph{formal} approach takes as many basis functions as instruments ($N=M$), piecewise constant on the region of sensitivity of each instrument, so that $J$ is square and (``if the basis functions are reasonable'') invertible, and solves $\bm x=J^{-1}\bm q^{\rm in}$ iteratively \ts{13}{10}{1:02:59}.
The iteration converges quickly because for at-the-money instruments the price is nearly proportional to the volatility: for the Black--Scholes forward price (\cref{ch:bs}), $C=F\,[2N(\sigma\sqrt T/2)-1]=F\,\sigma\sqrt{T}/\sqrt{2\pi}\,(1-\sigma^2T/24+\dots)$, so $\ln C\approx\ln\sigma+\text{const}$ and $q$ is nearly linear in the volatility
increments (my derivation of the speaker's remark; the calibrated instruments are swaptions in an HJM Monte Carlo, for which the proportionality is heuristic).

The outcome is an exact fit and a meaningless surface: ``Manhattan skyline, with the Hudson river'' \ts{13}{10}{1:04:01} (slide~33, ``Exact, but \dots\ meaningless''). Smoothing the piecewise-constant basis functions gives a better-looking surface but not a good one, and the instability
remains: a $1\%$ change in the price of the $5\text{y}\times10\text{y}$ swaption changes the adjustment of the surface by $4\%$ (slide~35: one tall isolated ``building with an antenna'') \ts{13}{10}{1:05:04}\ts{13}{10}{1:06:08}. This is what \cref{ch20:sec-noise} predicts: $J$ is square, so
$J^\dagger=J^{-1}$ and errors are amplified by up to $\kappa(J)$; the toy of \cref{ch20:cs-toy} shows the same phenomenon (1\% of noise, a factor $16$ on the surface).

\subsection{The regularised design}\label{ch20:sec-volsolve}
The lecture turns ill-posedness to advantage with four changes \ts{13}{10}{1:06:08}\ts{13}{10}{1:07:09} (slide~36).
\begin{itemize}
  \item \emph{Tolerance.} Each input has a bid--offer, so there is no point in calibrating exactly. The fit becomes a weighted least-squares term $\|W(\bm q-\bm q^{\rm in})\|^2$, $W$ diagonal with small weights on instruments of large bid--offer and large ones on important instruments. A
        small loss of accuracy buys a large gain of smoothness.
  \item \emph{Smooth basis functions.} Cubic B-splines (\cref{ch20:def-bspline}) in each direction, and their tensor products $B_a(t)B_b(s)$ as two-dimensional basis functions \ts{13}{10}{1:08:13}. Since B-splines are non-negative and sum to one (\cref{ch20:prop-bspline}), every combination is a
        reasonable surface, and its regularity is inherited from the basis (``every basis function makes sense, so any linear combination will be good enough'').
  \item \emph{More basis functions than instruments} ($N>M$): it is not known in advance which shapes are needed, and the number of unknowns is still tiny compared with the $28{,}920$ cells (with $24$ basis functions per direction, $576$ coefficients, so the matrix of the normal equations is $576\times576$, about $2.7$~MB) \ts{13}{10}{1:07:09}.
        The knots are denser at short forward times, where instruments are dense \ts{13}{10}{1:22:09}.
  \item \emph{Penalties.} The problem becomes \ts{13}{10}{1:09:19}
        \begin{equation}\label{ch20:eq-volpen}
          \|W(\bm q-\bm q^{\rm in})\|^2\to\min,\qquad \|L_1(\bm v-\bm v_0)\|^2\to\min,\qquad \|L_2\bm v\|^2\to\min ,
        \end{equation}
        with a penalty $L_1$ on the \emph{change} of the surface with respect to a prior and a penalty $L_2$ on the surface itself (slide~40). The gradient penalty is a matrix with rows $(\dots,1,-1,\dots)$ for every pair of neighbouring cells (slide~41, \ts{13}{10}{1:10:20}); on a rectangular
        grid of $n_t\times n_s$ nodes it is $L_1=\begin{psmallmatrix}D_t\otimes I\\ I\otimes D_s\end{psmallmatrix}$ with $D$ the one-dimensional first-difference matrices (on the triangle only the pairs of cells that both lie inside are kept).
\end{itemize}

Combining the three terms with weights $\lambda_1,\lambda_2$ gives (slide~42, \ts{13}{10}{1:10:20})
\begin{equation}\label{ch20:eq-volsol}
\begin{gathered}
  \bm x=\arg\min\Bigl\{\|W(J\bm x-\bm q^{\rm in})\|^2+\|\lambda_1L_1B\bm x\|^2+\|\lambda_2L_2\bm x\|^2\Bigr\}=K\,J\T W^2\bm q^{\rm in},\\
  K=\bigl(J\T W^2J+\lambda_1^2(L_1B)\T L_1B+\lambda_2^2L_2\T L_2\bigr)^{-1},
\end{gathered}
\end{equation}
which is \eqref{ch20:eq-gen} with the two penalty terms of \cref{ch20:sec-general} (setting the gradient of the objective to zero gives the normal equations; here $L_1B\bm x=\bm v-\bm v_0$ so the first penalty acts on the change of surface, and $L_2$ acts on the coefficients).
$K$ is invertible when the penalties together have no kernel in common with $WJ$; $\lambda_2>0$ with $L_2=I$ guarantees it. Two computational remarks from the Q\&A \ts{13}{10}{1:24:19}: only $J\T W^2\bm q^{\rm in}$ changes when market quotes tick, whereas the penalty matrices (and, if $J$ is
approximately constant, $J\T W^2J$) can be precomputed, so the calibration runs in real time on live data; and the linear algebra is done in the space of $N\ll28{,}920$ coefficients, never on the full surface. When the pricing map $\bm q(\bm x)$ is not linear, each Newton step solves a regularised problem of this form; the resulting damped iteration is a
Levenberg--Marquardt method (not discussed in class).

\begin{intuition}[Where in the chapter the pieces fit]
Formula \eqref{ch20:eq-volsol} is one object seen three ways. As \emph{statistics}: ridge regression (\cref{ch05:sec-shrink}) with a design matrix $WJ$, observation weights, and a prior that surfaces are smooth and close to $\bm v_0$. As \emph{linear algebra}: the SVD filter \eqref{ch20:eq-filter} on
the singular values of $WJ$ in the metric of the penalty. As \emph{engineering}: what a trader looks at is the residual: the surface is a smooth fit through the market, and what is left over is information.
\end{intuition}

Slide~43 shows the calibrated surface: smooth, with a few gentle peaks and a sharper one at short forward times. Asked about a peak, the speaker explains \ts{13}{10}{1:22:09}\ts{13}{10}{1:23:16} that the grid is uniform in quarters whereas the instruments are at
3~months, 6~months, 1, 2, 5, 10 years, so on a logarithmic axis the feature would not look sharp; the surface has more detail where there are more instruments, and no sharper features can be supported elsewhere.

\begin{finance}[Relative value: the calibrated surface as a benchmark]
Traders watch the surface as it moves with the live quotes \ts{13}{10}{1:24:19}. If an instrument trades away from the smooth surface implied by the others (something that ``is not typical'', a judgement that ``comes with years of experience'') it is a candidate mispricing: a spike is likely
to disappear soon, and one trades (a swaption, or a more exotic product with the appropriate exposure) against it. This is \emph{relative value analysis} and, in the speaker's words, how the desk makes money \ts{13}{10}{1:25:20}. It works only if the model is robust: a regulariser based on a fundamental
principle (smoothness) is more likely to be stable in the future than one estimated from a past window such as PCA \ts{13}{10}{1:26:24}. The same remarks apply to the swap curve.
\end{finance}

\begin{finance}[Interpolation, splines and their limits]
To a question on other fitting techniques the speaker answers that ``spline'' and ``interpolation'' are the same thing for him: one has limited inputs and wants the curve in between \ts{13}{10}{1:20:07}\ts{13}{10}{1:21:07}. The regularised fit goes beyond interpolation: it
does not pass through the data, which is the essential difference between \cref{ch20:sec-build} and \cref{ch20:sec-volsolve}.
\end{finance}

\paragraph{What regularisation buys, and costs} \ts{13}{10}{1:17:52}. It gives stability (indispensable for ill-conditioned problems) and a more realistic solution, at the expense of not fitting the data exactly and possibly of a \emph{biased} solution (for example a smoothness constraint makes the result smoother than reality). The bias
is controlled by \emph{what} is penalised, in the sense of the ``no bias on the kernel of $L$'' remark in \cref{ch20:sec-general}: open the textbook and one finds ``Tikhonov'' and penalises the amplitude, which would push the volatility towards zero; a flat surface is acceptable, so one penalises derivatives instead \ts{13}{10}{1:19:00}. The
slides' references are the Wikipedia pages on the HJM model, the yield curve, inverse problems, Tikhonov regularisation, the SVD, PCA and B-splines.

% ==================================================================
\begin{keyformulas}
\begin{itemize}
  \item Bond: $P=\sum_i\Lambda_iC_i$; yield $\Lambda_i=e^{-yt_i}$; duration $D=\sum_it_i\pi_i$, convexity $\mathcal C=\sum_it_i^2\pi_i$ \eqref{ch20:eq-yield}; bond spread $P=\sum_ie^{-st_i}\Lambda_iC_i$ \eqref{ch20:eq-spread}.
  \item Swap rate $C=\sum_ir_iw_i/\sum_iw_i$, $w_i=\delta_i\Lambda_i$ \eqref{ch20:eq-swaprate}.
  \item Cox--de Boor: $B_{i,k}=\frac{t-\tau_i}{\tau_{i+k}-\tau_i}B_{i,k-1}+\frac{\tau_{i+k+1}-t}{\tau_{i+k+1}-\tau_{i+1}}B_{i+1,k-1}$, $\sum_iB_{i,k}=1$.
  \item PCA hedge: $\bm x=R\T\bm r$, $R=P(H\T P)^{-1}$, $H\T R=I$ \eqref{ch20:eq-pcarisk}; smooth hedge: $R=(HH\T+\lambda^2(LJ)\T LJ)^{-1}H$ \eqref{ch20:eq-regRsol}.
  \item Inverse problem $\bm y=A\bm x+\bm\varepsilon$: $\bm x^\dagger=\sum_i\frac{\bm u_i\T\bm y}{s_i}\bm v_i$, error $\sum_i\frac{\bm u_i\T\bm\varepsilon}{s_i}\bm v_i$, $\kappa=s_1/s_r$ \eqref{ch20:eq-pinv}.
  \item Filters $\bm x_w=\sum_iw_i\frac{\bm u_i\T\bm y}{s_i}\bm v_i$: Tikhonov $w_i=\frac{s_i^2}{s_i^2+\lambda^2}$, $\bm x_\lambda=(A\T A+\lambda^2I)^{-1}A\T\bm y$; TSVD $w_i=\ind_{i\le k}$ \eqref{ch20:eq-tikh}; gain $\le1/(2\lambda)$; repricing $H_\lambda=\sum_iw_i\bm u_i\bm u_i\T\ne I$.
  \item General form $\min\|W(A\bm x-\bm y)\|^2+\lambda^2\|L\bm x\|^2$: $\bm x=(A\T W^2A+\lambda^2L\T L)^{-1}A\T W^2\bm y$; no bias on $\ker L$ \eqref{ch20:eq-gen}.
  \item MSE $=\sum_i[(1-w_i)^2\xi_i^2+\sigma^2w_i^2/s_i^2]+\sum_{i>r}\xi_i^2$; optimal $\lambda_i^2=\sigma^2/\xi_i^2$ \eqref{ch20:eq-mse}--\eqref{ch20:eq-wiener}.
  \item Choosing $\lambda$: discrepancy $\rho(\lambda)=\delta$; $\mathrm{GCV}=M\rho^2/(M-\operatorname{tr}H_\lambda)^2$.
  \item Surface calibration: $\bm v=\bm v_0+B\bm x$, $J=\partial\bm q/\partial\bm x$, $\bm x=KJ\T W^2\bm q^{\rm in}$ \eqref{ch20:eq-volsol}.
\end{itemize}
\end{keyformulas}

\begin{exercise}[(not from the course)]
\leavevmode
\begin{enumerate}[(a)]
  \item Show that the total number of cells of a triangular grid with $n$ quarters on a side is $n(n+1)/2$ and compute the memory (in GB, doubles) of the square matrix of the full problem for $n=240$.
  \item With a tensor-product cubic B-spline basis of $m$ functions per direction, how many coefficients, and what is the size of the matrix $K$ in \eqref{ch20:eq-volsol}?
  \item Explain why $J\T W^2J$ has rank at most the number of instruments, and why $K$ nevertheless exists.
\end{enumerate}
\end{exercise}
